\documentclass[letterpaper,conference]{IEEEtran}
\usepackage[T1]{fontenc}
\usepackage[utf8]{inputenc}
\usepackage{cite}
\usepackage{booktabs}
\usepackage{array}
\usepackage{balance}
\usepackage{graphicx}
\usepackage{url}
\usepackage[hidelinks]{hyperref}
\def\UrlBreaks{\do\/\do\-\do\_}
\providecommand{\StudyPrefix}{}
\input{\StudyPrefix generated/metrics.tex}
\input{\StudyPrefix generated/provenance.tex}
\title{Simulated Smart-Home Control with Small Language Models on SBCs: Task Success, Latency, and Replication}
\author{\IEEEauthorblockN{Hudson Moreira Guimarães dos Santos\IEEEauthorrefmark{1}, Pablo Cesar Calcina Ccori\IEEEauthorrefmark{1},\\
Laisa Caroline Costa De Biase\IEEEauthorrefmark{1}, Bruno Cavalcante de Souza Sanches\IEEEauthorrefmark{1},\\
Roseli de Deus Lopes\IEEEauthorrefmark{1}, Marcelo Knörich Zuffo\IEEEauthorrefmark{1}}
\IEEEauthorblockA{\IEEEauthorrefmark{1}University of São Paulo (USP), São Paulo, Brazil\\
Corresponding author: Hudson Moreira Guimarães dos Santos (hudson.moreira@usp.br)}}
\newcounter{examplecase}
\begin{document}
\maketitle
\begin{abstract}
Small language models can broaden the phrasing accepted by a domestic controller, but their usefulness also depends on correct decisions and response cost. We compare a finite-rule baseline and quantized Qwen2.5 models of 0.5B and 1.5B parameters on Labrador single-board computers. Hardware and inference are real; the home and all actions are simulated. Forty reserved Portuguese scenarios cover clear instructions, paraphrases, routines, ambiguity, and invalid requests, with two repetitions per condition. Mean scenario success was 52.5\% for rules, 37.5\% for 0.5B, and 15.0\% for 1.5B; request latency medians were 0.157, 36.848, and 94.663 s. All responses were structurally valid. The 0.5B configuration handled more paraphrases than rules, but both models struggled with rejection and clarification. Independent replicas on one, three, and six boards increased correct-response rates in fixed batches while retaining the same task-success counts. The study jointly evaluates decision quality and inference cost, distinguishing model proposals from executor interventions. These results identify a concrete deployment trade-off: the tested 0.5B model broadens paraphrase coverage at a substantial latency cost. The comparison concerns these configurations and this constructed workload.
\end{abstract}
\begin{IEEEkeywords}simulated smart home, small language models, SBC, task success, replication\end{IEEEkeywords}
\section{Introduction}
A domestic controller must interpret a request, preserve unrelated device states, and decide when to ask for information or reject an unsupported operation. These requirements make natural-language control more demanding than producing a valid action description. On a resource-constrained device, correctness must also justify the time and memory required to obtain a response.

Two research directions inform this problem: evaluating domestic instructions, including invalid ones, and characterizing language models under edge resource constraints. Their intersection motivates a focused deployment question: how do rules and two small-model configurations differ in correct execution of simulated domestic tasks, inference cost, and service of independent requests through replicas?

Local execution is motivated by retaining household requests on user-controlled hardware and keeping control available without an Internet connection. This study evaluates the resulting on-device inference cost; it does not measure privacy guarantees, compare cloud services, or establish that local processing is always preferable.

We address this question using real ARM inference and a two-room lighting simulator. The contributions are a comparison against a finite-rule reference on the same Labrador boards; an evaluation that distinguishes valid output, correct proposals, and executor intervention; and instrumented closed batches served by independent replicas. Each board runs a complete model. The replicas process separate requests rather than collaborating on a decision or partitioning a model.

\section{Related Work}
HomeBench~\cite{homebench} evaluates English instructions in a virtual home, covering valid, invalid, and mixed requests involving one or multiple devices. Its execution-success metric requires all requested device operations to succeed; an additional F1 metric describes operation-level performance. It compares language models and adaptation strategies rather than a matched Labrador deployment. Our smaller Portuguese workload focuses on missing-room clarification and unrelated-state preservation, with separate records for proposed, blocked, and applied actions. We use an independently constructed dataset and do not compare our success percentages with HomeBench scores.

Lu et al.~\cite{edge} characterize more than 60 small models using broad capability benchmarks and deployment-cost experiments. Capability evaluation uses a server, while cost evaluation includes smartphones and Jetson hardware, separating prefill and decoding and examining quantization. Our comparison instead measures task decisions and HTTP latency on the same CPU-only Labrador platform. It addresses a narrow control task rather than general capability rankings, connecting request cost to whether the resulting action is acceptable.

Qin et al.~\cite{empirical} examine how resource constraints affect personalized-model choices, including learning methods, data volume, compression, task difficulty, and available learning time. Their focus on customization differs from our fixed pretrained configurations: we perform inference without training or personalization. Together, these studies motivate joint evaluation of task requirements and deployment cost. Our addition is a traceable comparison of rules, small models, and independent replication for a Portuguese simulated-control workload. The comparisons are methodological; differences in tasks, models, language, and hardware preclude direct numerical rankings across studies.

\section{Experimental Method}
\subsection{Simulator and execution policy}
A Python coordinator creates a fresh state for each request, calls a resident HTTP service, validates its JSON, applies permitted changes, and independently evaluates the result. No domestic device is connected. The two rooms are identified by the technical values \texttt{sala} (living room) and \texttt{quarto} (bedroom). Lights take \texttt{on}/\texttt{off} values. The contract admits \texttt{set\_light}, \texttt{set\_lights}, \texttt{ask\_clarification}, and \texttt{no\_action}; a two-step plan is atomic, so one permission failure blocks the whole plan.

The executor sees only the request, visible state, and permissions, without access to evaluator answers or scenario categories. A conservative policy blocks lighting proposals if neither a known room name appears as a whole word in the request nor a valid location is supplied. Unknown rooms and unauthorized values are also blocked. This policy checks permission, not task correctness: a permitted proposal can still change the wrong light or misinterpret an invalid instruction. Proposals, intervention, and outcomes are recorded separately.

\subsection{Hardware and configurations}
A notebook runs the Python experiment coordinator and sends text requests over HTTP/IPv4 through an Ethernet LAN (192.168.50.0/24) to each board. The notebook and all boards use Ethernet for this experiment. The coordinator is a client and evaluator, not a model shard. Each Labrador hosts a complete model; replicas do not exchange intermediate activations or collaborate on one decision. Inputs in this campaign are exclusively text, with no audio, camera, or physical actuator in the measurement path.

Six Labrador boards provide four ARM64 CPU cores, Debian 12, and approximately 1.89 GiB of OS-visible RAM each. Models use \texttt{llama-server}~\cite{llamaserver} build 9584/\texttt{e25a32e98}, CPU execution, four compute and four batch threads, a 512-token context, and one slot per model service. Both model services stay resident on each board, with at most one inference active across them. Qwen2.5-0.5B-Instruct and Qwen2.5-1.5B-Instruct~\cite{qwen05,qwen15} use Q5\_K\_M GGUF files whose hashes match the official Qwen files~\cite{qwen05gguf,qwen15gguf}. Runtime and library hashes were checked across boards; full provenance is in the artifact.

Rules run as a resident stateless Python HTTP service on the same boards. The finite grammar recognizes fixed Portuguese verbs for turning lights on or off, explicit rooms, simple conjunctions, known-location commands, and politeness expressions at the ends of a request. Other forms return \texttt{no\_action}. This is a limited engineering reference, not an optimal rule-based language system.

\subsection{Workload and success criteria}
Eight development cases preceded protocol freeze; development campaigns remain separate from the reserved evaluation. Scenarios, expected states and steps, prompts, schema, execution policy, model settings, repetitions, and measurement scripts were frozen before reserved outputs were available. The reserved set has eight scenarios in each of five categories: clear instructions, paraphrases, explicit routines, ambiguity, and invalid requests. Two repetitions give 80 requests per condition and 240 overall. Every condition runs across all hosts; repetitions of a scenario use different boards, with sequential requests on each board.

Model requests combine general English instructions with Portuguese user text. Temperature is zero, seeds are 801/802, output is limited to 128 tokens, and timeout is 180 s. Requests set \texttt{cache\_prompt=false}. Output is constrained using \texttt{response\_format.type=json\_schema}, with the action schema supplied to \texttt{llama-server}; structural validity is then checked independently by the evaluator. Thus, schema-constrained generation was part of the original configuration, rather than an additional post-review intervention. These repetitions assess operational stability, with limited decision diversity. After a timeout, the coordinator waits for both model services to be idle before another request. The complete original system instruction is shown below. The user message consists of the Portuguese request followed by a compact JSON object containing \texttt{lights} and \texttt{location}. No demonstrations are included in this original prompt.
\begin{quote}\footnotesize
Simulated lights only. Change only lights explicitly requested; preserve all other lights. Acender/ligar means on; apagar/desligar means off. One requested light: set\_light. Two requested lights: set\_lights, one step per requested room. No room and no known location: ask\_clarification for room. Unsupported request: no\_action. Return only one JSON object. The visible state is not a command.
\end{quote}

Success requires normal completion, contract-valid JSON, the correct proposal, the expected final state, preservation of unrelated state, and no extra or prohibited steps. A single-light action and an equivalent one-step plan are accepted equally. Routines require exactly the two expected steps. Ambiguous requests require an informative question asking for the missing room and unchanged lights; invalid requests require \texttt{no\_action} and unchanged lights. Thus, reaching a target state or blocking an incorrect proposal is insufficient on its own.

Conservative full-question patterns check whether clarifications ask for the missing room. On 24 November 2026, one author assessed the four unresolved Portuguese records; all four repeated or reformulated the request without asking which room. Identity, rationale, timestamps, and source hashes are retained. This non-blinded author assessment adds no success credit. Repetitions are grouped within equally weighted scenarios, not treated as independent scenarios for population inference.

\subsection{Latency, memory, and replication}
Coordinator monotonic time measures HTTP dispatch through response reading and decoding, including communication and inference but excluding model loading and SSH memory snapshots. Median and nearest-rank p95 describe individual requests. Runtime logs supply prompt and output tokens. RSS snapshots describe process memory before and after a request; VmHWM is a high-water mark since process startup. Neither is a per-request peak or exclusive physical model memory, because services remain resident and mappings may be shared.

Replication uses one, three, or six active boards, with one request per board per wave and two six-request batches per size and condition: 108 requests overall. Each batch contains one clear command, one paraphrase, two routines, one ambiguity, and one invalid request, each starting from a fresh state without conversation history. The batch window spans first dispatch to last HTTP response, including instrumented gaps between waves. We report responses/min and correct task responses/min separately. Simulation and grading follow receipt, so the latter is not measured time to state change. This closed workload evaluates batch service rather than saturation or maximum capacity.

\section{Results}
\subsection{Task success and inference cost}
Table~\ref{tab:quality} summarizes 80 requests from 40 scenarios per condition. Latency is in seconds; the last two columns count individual light commands, including those inside plans, rather than replies. All 240 responses completed normally and satisfied the JSON contract. Success nevertheless differed substantially: rules completed 42/80 attempts, 0.5B completed 30/80, and 1.5B completed 12/80.
\begin{table}[t]\caption{Task success, latency, and light commands.}
\centering\footnotesize\setlength{\tabcolsep}{4pt}\begin{tabular}{lrrrrr}\toprule
Condition & Success & Median & p95 & Undue & Blocked\\\midrule
\input{\StudyPrefix generated/quality_rows.tex}
\end{tabular}\label{tab:quality}\end{table}

The category comparison in Table~\ref{tab:categories} identifies where the approaches differed. Each cell covers eight scenarios with two repetitions, or 16 attempts; these are not 16 independent scenarios. Rules completed every clear and invalid request but none of the paraphrases. The 0.5B configuration completed 14/16 paraphrases, showing broader coverage of these formulations, yet rejected none of the invalid requests correctly. Neither model successfully handled ambiguity under the recorded criteria.
\begin{table}[t]\caption{Task success by request category.}
\centering\footnotesize\begin{tabular}{lrrr}\toprule
Category & Rules & 0.5B & 1.5B\\\midrule
\input{\StudyPrefix failure_analysis/category_rows_en.tex}
\end{tabular}\label{tab:categories}\end{table}

The author assessment confirmed all four previously unresolved 0.5B clarification failures. Aggregate success remains 30/80 (37.5\%) for 0.5B; no success was awarded through manual assessment. All credited successes remain automatic.

The 0.5B configuration also had lower request latency than 1.5B: medians were 36.848 and 94.663 s, versus 0.157 s for rules. Median post-request RSS was 571.6 and 1183.7 MiB for the models. Both used 124--138 prompt tokens; output ranges were 16--40 and 12--75 tokens, respectively. Recorded prompt-cache counts were zero. These measurements characterize the complete configurations rather than isolating the effect of parameter count.

\subsection{Decision failures and representative cases}
Table~\ref{tab:failures} groups existing evaluator checks without changing outcomes. Every entry gives failed/observed attempts. Step and target-state checks apply to 48 clear, paraphrase, and routine attempts; preservation applies to 64 attempts, excluding routines. Rows overlap and cannot be summed. Prohibition checks exist in all 80 attempts, although only 32 ambiguity/invalid attempts have nonempty prohibitions. Clarification and rejection rows evaluate proposals of those types; no clarification proposal is not evidence of clarification ability.
\begin{table}[t]\caption{Failed requirements among observed attempts.}
\centering\footnotesize\begin{tabular}{lrrr}\toprule
Requirement & Rules & 0.5B & 1.5B\\\midrule
\input{\StudyPrefix editorial_20261029/failure_rows_en.tex}
\end{tabular}\label{tab:failures}\end{table}

Rules preserved unrelated state but failed to act on valid requests outside their grammar. The models exhibited a different profile: incorrect step sets, prohibited proposals, and changes to unrelated lights. In particular, 1.5B failed step equality in 38/48 applicable attempts but the target-state check in 22/48, illustrating why the requested light alone cannot establish task completion. At command level, 0.5B proposed 36 undue changes, of which 16 were blocked and 20 simulated; 1.5B proposed 76, of which 22 were blocked and 54 simulated. These command counts differ from blocked-attempt counts because a reply may contain multiple commands and an entire plan is blocked atomically.

Table~\ref{tab:cases} presents four recorded cases. Requests and response text are English translations for presentation; the experiments used Portuguese. Technical JSON keys and values are retained, and action descriptions abbreviate the recorded JSON. Original utterances, complete outputs, states, and the mapping from case labels to attempt directories are preserved in the supplement. The inadequate-clarification count includes the four failures confirmed by the author assessment.
\input{\StudyPrefix editorial_20261029/examples_en.tex}

Case~\ref{case:paraphrase} shows a valid paraphrase rejected by rules. Case~\ref{case:room} separates an incorrect room assumption from its blocked execution. Case~\ref{case:clarification} illustrates a question-field response that fails to ask for the missing information. Case~\ref{case:preservation} reaches the requested light state but changes the other room, failing preservation and exact-step criteria. These cases explain the evaluation distinctions without attributing unmeasured internal causes to the models.

\subsection{Independent replicas}
The closed-batch experiment measures service of separate requests by complete model copies. With one, three, and six boards, Qwen 0.5B produced mean correct-response rates of 0.871, 2.883, and 6.393/min; Qwen 1.5B produced 0.085, 0.289, and 0.661/min. Each six-request batch retained three successes for 0.5B, one for 1.5B, and five for rules. These are descriptive rates from two batches per size and condition, not sustained capacity estimates.

The six-replica rule batches lasted only about 9--11 ms, making their derived per-minute rate especially sensitive to instrumentation and the absence of between-wave gaps. We retain those raw windows and rates in the artifact and omit that rate from the main comparison. Parallel copies increased concurrent service; they did not correct interpretation errors. This experiment is a deployment baseline, not a new parallelization algorithm.

\section{Discussion}
\subsection{Choosing an approach for the task}
Rules offered the lowest cost and highest aggregate success, but failed every paraphrase. The 0.5B configuration covered more paraphrases while failing rejection and preservation checks. Controller choice therefore depends on supported requests, acceptable errors, latency, and memory. A rules-first controller with language-model fallback requires separate evaluation: rule rejection cannot distinguish an unfamiliar valid paraphrase from an invalid request. The present results do not validate such a hybrid.

The distinction between syntax and execution complements HomeBench's emphasis on complete instruction execution~\cite{homebench}. Our missing-room and preservation cases further show how a blocked proposal or a correct target light can mask an unsuccessful decision. Their value is diagnostic, not evidence that our smaller workload is more comprehensive than HomeBench.

The comparison extends the deployment perspective of Lu et al.~\cite{edge}: footprint and valid JSON alone are insufficient for control. The 1.5B configuration failed all 16 clear-command attempts, sometimes changing an unrelated light despite a correct target. This warrants checking prompt and runtime configuration, rather than attributing failure to parameter count. Qin et al.~\cite{empirical} likewise emphasize task and resource constraints, but their customization setting cannot explain this ranking.

A separate exploratory follow-up reused the eight development cases with Qwen 1.5B, one repetition, and the original runtime, schema, seed 801, and decoding settings. Original English instructions, Portuguese instructions, and Portuguese instructions plus two examples yielded 1/8, 3/8, and 4/8 successes, respectively; medians were 78.5, 93.2, and 124.0 s. All 24 outputs were structurally valid. These are development observations, not revised reserved-set scores. Five recovered hosts were used; after one became unreachable, its three remaining requests moved to a host with matching model/runtime hashes. Development reuse, one repetition, and host reassignment preclude causal or population claims; the results motivate an independently frozen prompt comparison.

The original 36.848 and 94.663 s medians remain a practical obstacle to interactive control. Smaller or distilled models, shorter prompts, caching, and acceleration require separate configuration and latency evaluation; no such speedup is established here.

Replication addresses a separate deployment decision. Parallel copies increased correct-response rates for independent requests, with similar individual medians across replica counts and unchanged success counts in each batch. They did not improve interpretation of a request. Evaluating sustained service would require an arrival process, queue measurements, and longer runs beyond these instrumented batches.

\subsection{Validity and scope}
The 40 authored utterance/state combinations cover two lights in Portuguese, not independent homes or a population of users. Several express the same underlying action. Requested lights start in the opposite state, leaving idempotent requests and other state distributions untested; routine success alone cannot distinguish language understanding from state inversion. The finite grammar is a limited baseline, and one prompt, schema, and quantization per model do not isolate architecture or size effects.

Host interleaving reduces confounding without eliminating thermal, memory, or network variation. Resident services and shared mappings limit memory interpretation; page caches were not cleared, and energy and per-request peaks were not measured. Few closed batches, including instrumentation costs, limit service-rate generalization. The four clarification judgments have a single-author assessment, without blinded or independent multi-rater validation. Simulation provides no evidence of real-home safety or usability. These limits define the scope of the configuration comparison rather than changing its observed outcomes.

\section{Conclusion}
For the tested simulated domestic tasks, rules achieved the highest aggregate success at the lowest response cost, while Qwen 0.5B handled more paraphrases than the finite grammar and outperformed the tested 1.5B configuration in both success and latency. Valid structured output did not guarantee correct decisions: rejection, clarification, and preservation required separate evaluation. Independent replicas increased correct-response rates in closed batches while leaving task-success counts unchanged.

Selecting a controller therefore requires matching supported requests and decision criteria to the available inference budget. A concrete next step is a new, independently reviewed evaluation with user-authored phrasing, idempotent commands, and varied initial states, retaining the distinction between proposed actions, executor intervention, and task success.

\section*{Evidence Availability and AI-use Disclosure}
The frozen protocol, 240 quality and 108 replication attempts, review ledger, and reconstruction scripts are available at an immutable snapshot~\cite{artifact}. The \href{https://github.com/hudsonmoreira-spec/labrador-distributed-llm-benchmarks/tree/f2f566b179fcf3893099d7d1b12b888455e531ae/experiments/smarthome_agents/analysis/failure-review-20261028}{failure analysis} and \href{https://github.com/hudsonmoreira-spec/labrador-distributed-llm-benchmarks/tree/c0650b876b5e1a9169782d9cc0819c078b13fbe9/paper/smarthome_agents/study_v1/editorial_20261029}{translation notes} preserve example provenance. Separate revision artifacts retain diagnostic plans, raw exchanges, host changes, instrumentation errors, and author assessments with source hashes. Hardware replication requires the recorded runtime and models.

No external funding was received. OpenAI Codex assisted code, analysis, figures, evidence checks, and manuscript revision. The author reports comparing his own clarification assessments with assistant explanations and using suggested wording for rationales; this is not independent multi-rater validation. Authors remain responsible for verification and the final text under IEEE guidance~\cite{ieeeai}.
\bibliographystyle{IEEEtran}
\bibliography{\StudyPrefix references}
\end{document}
